\documentclass[11pt,a4paper,sans]{moderncv}
\moderncvstyle{classic}
\moderncvcolor{blue}
\usepackage[utf8]{inputenc}
\firstname{Zein}
\familyname{ELAJAMY}
\title{Ingénieur Simulation Numérique | Éléments Finis \& Flambage Dynamique | R\&D Industrielle}
\phone{+33 7 54 45 47 42}
\email{zeinfrance4@gmail.com}

\begin{document}
\makecvtitle

\section{Résumé}
Ingénieur R\&D spécialisé en simulation numérique et performance énergétique industrielle. Éléments finis (Abaqus/Metafor), CFD (Ansys Fluent), jumeaux numériques et outils Python d'analyse à grande échelle. Expérience terrain chez ArcelorMittal : modèle thermique validé à 95\%, trajectoire de décarbonation orientée par audit Python/Plotly. Flambage dynamique de tôles minces (PFE en cours).

\section{Expériences}
\cventry{Avril 2026 - Août 2026}{Projet de Fin d'Études : Flambage Dynamique des Tôles Minces Laminées}{ArcelorMittal R\&D}{None}{}{
\begin{itemize}
\item Modélisation de cas tests sous Abaqus et Metafor. Réalisation d'études de sensibilité numérique (maillage, type d'éléments, conditions de symétrie). Optimisation du couple maillage/pas de temps et comparaison des solveurs Implicite vs Explicite. Analyse comparative des performances et de la stabilité entre Metafor et Abaqus.
\end{itemize}
}
\cventry{Année 1 (2023-2024)}{Jumeau Numérique Thermique}{ArcelorMittal R\&D}{None}{}{
\begin{itemize}
\item Développement d'un modèle thermique hors-ligne en Python pour simuler le refroidissement industriel. Utilisation de corrélations physiques et intégration de données géométriques. Structuration du dossier logiciel pour la reproductibilité.
\end{itemize}
}
\cventry{2023}{Stagiaire Recherche - Thermodynamique des Alliages}{Institut Jean Lamour (IJL)}{None}{}{
\begin{itemize}
\item Utilisation de logiciels spécialisés (ThermoCalc) pour le calcul de diagrammes de phases. Analyse et interprétation des résultats expérimentaux.
\end{itemize}
}
\cventry{Semestre 9}{Comportement Thermomécanique d’un Composite Stratifié}{ENSEM Nancy}{None}{}{
\begin{itemize}
\item Approche analytique et simulation Abaqus/CAE. Analyse de contraintes résiduelles.
\end{itemize}
}
\cventry{Année 2 (2024-2025)}{Jumeau Numérique Énergétique}{ArcelorMittal R\&D}{None}{}{
\begin{itemize}
\item Développement d'un module de calcul en Python pour les bilans énergétiques. Implémentation d'une logique de calcul multi-niveaux et inversion mathématique de corrélations physiques. Création de dashboards interactifs pour la visualisation des données.
\end{itemize}
}

\end{document}
